% !TEX root = ../main.tex
\section{Clasificación de los datos según su Tipo de Dato Abstracto}

Una vez identificados los datos que forman parte del expediente de cada
estudiante, es necesario determinar la naturaleza de la información que
representa cada uno. Esta clasificación permite definir posteriormente
la forma en que los datos serán almacenados y manipulados dentro del
programa.

Para este diseño se consideran tipos como entero, real, cadena de
caracteres, dato enumerado y fecha. La clasificación lógica de cada dato
se presenta en la siguiente tabla.

\begin{longtable}{p{3.8cm} p{3.2cm} p{7cm}}
\toprule
Dato & Tipo de dato & Justificación \\
\midrule
\endhead

Número de carné &
Entero &
Representa el identificador numérico asignado por la universidad a cada
estudiante. Se utilizará principalmente como llave de búsqueda y
referencia del registro. \\

Número de identificación &
Cadena de caracteres &
Aunque está compuesto principalmente por dígitos, funciona como un
identificador y puede variar en longitud o incluir caracteres según el
tipo de documento. Por esta razón se trata como información textual. \\

Nombre &
Cadena de caracteres &
Está compuesto por una secuencia de caracteres alfabéticos y forma parte
de la información personal del estudiante. \\

Apellidos &
Cadena de caracteres &
Corresponde a información textual. En el diseño posterior podrá
separarse en diferentes campos para facilitar búsquedas y ordenamientos. \\

Fecha de nacimiento &
Fecha &
Está formada por varios valores relacionados entre sí, específicamente
día, mes y año. Por esta razón resulta conveniente representarla mediante
un tipo compuesto destinado al manejo de fechas. \\

Sexo &
Enumerado &
Corresponde a un conjunto limitado de valores posibles, por lo que puede
representarse mediante un dominio cerrado en lugar de utilizar texto
libre. \\

Carrera &
Entero de referencia &
Puede almacenarse mediante un código numérico que identifique una carrera
dentro de un catálogo, evitando repetir el nombre completo de la carrera
en cada registro. \\

Nivel académico &
Enumerado &
Representa una categoría dentro de un conjunto limitado de niveles
académicos definidos por el sistema. \\

Correo electrónico &
Cadena de caracteres &
Puede contener letras, números y símbolos, por lo que debe almacenarse
como una secuencia de caracteres. \\

Número telefónico &
Cadena de caracteres &
Aunque contiene dígitos, no representa una cantidad sobre la cual deban
realizarse operaciones matemáticas. Se utiliza únicamente como dato de
contacto. \\

Dirección &
Cadena de caracteres &
Contiene texto de longitud variable que puede incluir letras, números,
espacios y otros caracteres. \\

Cantidad de créditos aprobados &
Entero &
Representa una cantidad discreta de créditos acumulados por el estudiante
y puede utilizarse en operaciones de comparación y actualización. \\

Promedio general &
Real &
Puede contener una parte decimal, por lo que requiere un tipo numérico
capaz de representar valores fraccionarios. \\

Estado del estudiante &
Enumerado &
Representa un conjunto limitado de condiciones posibles del estudiante,
como activo o inactivo. \\

\bottomrule
\end{longtable}

Esta clasificación distingue entre los datos que representan cantidades
numéricas y aquellos que, aunque puedan contener dígitos, funcionan
únicamente como identificadores. Por ejemplo, el número telefónico y el
número de identificación no se utilizan para realizar cálculos, mientras
que la cantidad de créditos aprobados y el promedio general sí forman
parte de operaciones académicas.

También se identifican datos cuyo dominio está limitado a un conjunto
determinado de valores, como el sexo, el nivel académico y el estado del
estudiante. Estos pueden tratarse como tipos enumerados para evitar el
almacenamiento de valores arbitrarios.

La clasificación anterior establece la representación lógica de los
datos. En las siguientes secciones estos tipos se traducen a estructuras
concretas del lenguaje C++, respetando las restricciones del proyecto y
las necesidades de almacenamiento del sistema.